\documentclass[a4paper,10 pt,titlepage,twoside]{book}

\usepackage{color, graphicx}
%\usepackage[luatex]{color, graphicx}

\usepackage{mathrsfs}
\usepackage{array,epsfig}
\usepackage{amsmath}
\usepackage{eufrak}
\usepackage{amsfonts}
\usepackage{mathbbol}
\usepackage{amssymb}
\usepackage{amsxtra}
\usepackage{amsthm}
\usepackage{stmaryrd}
\usepackage{tikz-cd}
\usepackage{enumerate}
\usepackage{tikz}
\usepackage{lmodern}
\usepackage[english]{babel}
\usepackage{geometry}
\geometry{a4paper,top=3cm,bottom=3cm,left=3.5cm,right=3.5cm,heightrounded,bindingoffset=5mm}
\usepackage{booktabs}
\usepackage{listings}
\usepackage{emptypage}
\usepackage{hyperref}
\usepackage{cleveref}
\usepackage{enumerate}
\usepackage{indentfirst}
%\usepackage{MnSymbol}

\usepackage{lipsum}

\newcommand{\Z}{\mathbb{Z}}
\newcommand{\Q}{\mathbb{Q}}
\newcommand{\R}{\mathbb{R}}
\newcommand{\C}{\mathbb{C}}
\newcommand{\K}{\mathbb{K}}
\newcommand{\CP}{\mathbb{CP}}
\newcommand{\N}{\mathbb{N}}
\newcommand{\LC}{\mathcal{L}}
\newcommand{\support}{\mathrm{support}}
\newcommand{\Oset}{\varnothing}
\DeclareMathOperator{\id}{id}

\numberwithin{equation}{chapter}
\renewcommand{\theequation}{\thechapter.\arabic{equation}}
%per gli spazi:
\usepackage{setspace}
\singlespacing

\providecommand\phantomsection{}




%\usepackage{stmaryrd} 

\theoremstyle{plain}
\newtheorem{thm}{Theorem}[chapter] % reset theorem numbering for each chapter

\theoremstyle{definition}
\newtheorem{defn}[thm]{Definition} % definition numbers are dependent on theorem numbers
\newtheorem{ex}[thm]{Example} % same for example numbers
\newtheorem{rmk}[thm]{Remark}
\newtheorem{notn}[thm]{Notation}
\newtheorem{lem}[thm]{Lemma}
\newtheorem{cor}[thm]{Corollary}
\newtheorem{prop}[thm]{Proposition}

%Here I define some theorem styles and shortcut commands for symbols I use often





%Pagination stuff.
\setlength{\topmargin}{-.3 in}
\setlength{\oddsidemargin}{0in}
\setlength{\evensidemargin}{0in}
\setlength{\textheight}{9.in}
\setlength{\textwidth}{6.5in}
\pagestyle{empty}


\DeclareMathOperator{\obj}{Obj}
\DeclareMathOperator{\tr}{tr}
\DeclareMathOperator{\mor}{Mor}
\DeclareMathOperator{\iso}{Iso}
\DeclareMathOperator{\ve}{vect}
%\DeclareMathOperator{\diff}{diffmfd}
\DeclareMathOperator{\g}{\mathfrak{g}}
\DeclareMathOperator{\1}{\mathbb{1}}

%\DeclareMathOperator{\del}{\partial\!}
\newcommand{\del}{\partial}
\DeclareMathOperator{\grad}{grad}
\DeclareMathOperator{\ind}{ind}
\DeclareMathOperator{\ch}{ch}
\DeclareMathOperator{\td}{Td}
\DeclareMathOperator{\pf}{pf}   
\DeclareMathOperator{\cl}{Cl}  
\DeclareMathOperator{\htimes}{\hat{\otimes}}   
\DeclareRobustCommand{\rchi}{{\mathpalette\irchi\relax}}
\newcommand{\irchi}[2]{\raisebox{\depth}{$#1\chi$}} % inner command, used by \rchi

\newcommand{\vect}[1]{\ensuremath{\begin{pmatrix}#1\end{pmatrix}}}

% For Spin Geometry
\DeclareMathOperator{\diff}{\operatorname{Diff}}

% Lineare Algebra %
\DeclareMathOperator{\rg}{rg}     % Rang
\DeclareMathOperator{\diag}{diag} % Diagonalmatrix
\DeclareMathOperator{\Mat}{M}   % Matritzenraum
\DeclareMathOperator{\Eig}{Eig}   % Eigenräume
\DeclareMathOperator{\End}{End} 
\DeclareMathOperator{\Span}{Span} % lineare Huelle
\DeclareMathOperator{\adj}{adj}   % Adjungierte
\DeclareMathOperator{\Ann}{Ann}   % Annihilator
\DeclareMathOperator{\Aut}{Aut}   % Automorphismen Gruppe
\DeclareMathOperator{\Sp}{Sp}     % Symplektische Gruppe
\DeclareMathOperator{\GL}{GL}     % Allgemeine Lineare Gruppe

\DeclareMathOperator{\Tr}{Tr}     % Spur
\DeclareMathOperator{\Ker}{Ker}   % Kern 
\DeclareMathOperator{\Bild}{Im}   % Bild
\DeclareMathOperator{\Hom}{Hom}   % Homomorphismen
\DeclareMathOperator{\sgn}{sgn}   % Signum
\DeclareMathOperator{\Deg}{Deg}   % Degree
\DeclareMathOperator{\vol}{vol}   % Signum


%\usepackage[backend=bibtex,
%style=alphabetic,
%sorting=ynt]{biblatex}
%\addbibresource{Lit.bib}

\setlength\parindent{24pt}


\newcommand{\bigzero}{\mbox{\normalfont\Large\bfseries 0}}
\newcommand{\rvline}{\hspace*{-\arraycolsep}\vline\hspace*{-\arraycolsep}}

\usepackage{wasysym}
\newcommand{\heart}{\gamma}
\renewcommand\epsilon{\varepsilon}
\renewcommand\phi{\varphi}


\usepackage{fontspec}
\newfontfamily\TahomaFont[
Scale=1,Ligatures={TeX},
BoldFont={Tahoma Bold}, 
ItalicFont={Open Sans Italic},
]{Tahoma}

\usepackage[absolute]{textpos}

\begin{document}

\thispagestyle{empty}

\begingroup
\TahomaFont

\begin{textblock*}{4cm}(8.65cm,1.03cm)
	\centerline {\includegraphics[width=3.67cm]{LogoUniToConvertito.png}}
\end{textblock*}


\begin{textblock*}{21cm}(0cm,8,98cm)
	\fontsize{18}{22}\selectfont
	\centerline {\textbf{ Universit\`a degli Studi di Torino}}
\end{textblock*}
\begin{textblock*}{21cm}(0cm,9.97cm)
	\fontsize{18}{22}\selectfont
	\centerline {\textit{ Corso di Laurea in Fisica}}
\end{textblock*}


\begin{textblock*}{21cm}(0cm,13.66cm)
	\fontsize{20}{24}\selectfont
	\centerline  {\textbf{Testo di sopra}}
\end{textblock*}
\begin{textblock*}{21cm}(0cm,14.46cm)
	\fontsize{18}{22}\selectfont
	\centerline{\Large {Tesi di Laurea}}
\end{textblock*}



\fontsize{14}{17}\selectfont

\begin{textblock*}{8cm}(3.04cm,18.26cm)
	\noindent 
	\textbf{Relatrice}
\end{textblock*}
\begin{textblock*}{8cm}(3.04cm,18.85cm)
	\noindent 
	Maria Luisa Frau
\end{textblock*}

%\begin{textblock*}{8cm}(3.04cm,20.64cm)
%	\noindent 
%	\textbf{Correlatore/Correlatrice}
%\end{textblock*}
%\begin{textblock*}{8cm}(3.04cm,21.23cm)
%	\noindent
%	Cognome Nome
%\end{textblock*}

\begin{textblock*}{8cm}(12.33cm,22.95cm)
	\noindent
	\textbf{Candidato/a}
\end{textblock*}
\begin{textblock*}{8cm}(12.33cm,23.55cm)
	\noindent	\textbf{Rocchietti March Paolo}
\end{textblock*}
\begin{textblock*}{8cm}(12.33cm,24.14cm)
	\noindent	Matricola 1093935
\end{textblock*}

\begin{textblock*}{21cm}(0cm,27.34cm)
	\centerline{Anno Accademico 2025/2026}
\end{textblock*}

\endgroup


\newpage
$ $

\TPoptions{absolute=false}
\pagestyle{plain}
\frontmatter

%\newpage

\tableofcontents



\newpage
$ $

\chapter*{Introduction}
\addcontentsline{toc}{chapter}{Introduction}



\newpage
$ $

\mainmatter
\chapter{Title of the first chapter}
\label{chap:1}

\section{Title of the first section of the first chapter}
\label{sec:11}
\begin{defn}
    A definition
\end{defn}


\begin{ex}
    An example
\end{ex}


A paragraph ...

\begin{prop}
\label{corr:3}

A proposition

\end{prop}



\section{Title of the second section of the first chapter}
\label{sec:12}
\begin{rmk}
A remark
\end{rmk}

\noindent
A list of formulae

\begin{enumerate}
    \item $t_{ii}(p)=p$ if $p \in U_i$;
    \item $t_{ij}(p) = t_{ji}(p)^{-1}$ for any $p \in  U_i \cap U_j$;
    \item $t_{ij}(p)\cdot t_{jk}(p) = t_{ik}(p)$ for any $p\in U_i \cap U_j \cap U_k$.
\end{enumerate}
A diagram
\[
\begin{tikzcd}
    f^*E \arrow{d}{} & E \arrow{d}{}\\
    M' \arrow{r}{} & M
\end{tikzcd}
\]
A numbered displayed formula
\begin{equation}
\label{ps:2}
    (f^*E)_p= E_{f(p)}.
\end{equation}

An unnumbered displayed formula
\[
	a=b \wedge b=c \Rightarrow a=c
\]

Notation explanation
\begin{notn}
If $k$ is clear form the context, we will simply write $\sigma(P,\xi)$.
\end{notn}

A lemma

\begin{lem}
\label{corr:15}
Statement of the lemma
\end{lem}

\chapter{Title of the second chapter}
\label{chap:2}
\section{Title of the first section of the second chapter}
\label{sec:21}

A corollary

\begin{cor}
\label{cor:2}
    Statement of the corollary
    
    \end{cor}

    
\section{Title of the second section of the second chapter}
\label{sec:22}

A theorem 
\begin{thm}
    Statement of the theorem, including an unnumbered displayed formula
    \begin{displaymath}
        HP_k\left(\sum_{j=1}^n g(X_j)\right)= \alpha_k(\sigma_{n,1}, \dots, \sigma_{n,k}).
    \end{displaymath}
\end{thm}

\newpage

An inline formula at the top of the page: $\sin^2 \alpha+\cos^\alpha=1$, for any $\alpha\in\R$.
    
    
\chapter{Title of the third chapter}
A quotation \cite{APS1}

A reference to a numbered formula \eqref{ps:2}

A reference to the first section of the first chapter \ref{sec:11}

\chapter*{Ringraziamenti}

Ringraziamenti

\phantom{p. 1}
\clearpage
\thispagestyle{empty}
\phantom{p. 2}
\clearpage
\addcontentsline{toc}{chapter}{References}

\begin{thebibliography}{9}
\bibitem{APS1}
M. F. Atiyah, V. K. Patodi, I. M. Singer (1975), \emph{Spectral asymmetry and Riemannian geometry. I}, Math.Proc. Cambridge Philos. Soc. \textbf{77}, 43-69.

\bibitem{APS2}
M. F. Atiyah, V. K. Patodi, I. M. Singer (1975), \emph{Spectral asymmetry and Riemannian geometry. II}, Math.Proc. Cambridge Philos. Soc. \textbf{78}, no. 3, 405-432.

\bibitem{APS3}
M. F. Atiyah, V. K. Patodi, I. M. Singer (1976), \emph{Spectral asymmetry and Riemannian geometry. III}, Math.Proc. Cambridge Philos. Soc. \textbf{79}, no. 1, 71-99.

\end{thebibliography}


$\R$
\end{document}
